%--------------------
% Packages
% -------------------
\documentclass[11pt,a4paper]{article}
\usepackage[utf8x]{inputenc}
\usepackage[T1]{fontenc}
%\usepackage{gentium}
\usepackage{mathptmx} % Use Times Font

\usepackage[pdftex]{graphicx} % Required for including pictures
\usepackage[swedish]{babel} % Swedish translations
\usepackage[pdftex,linkcolor=black,pdfborder={0 0 0}]{hyperref} % Format links for pdf
\usepackage{calc} % To reset the counter in the document after title page
\usepackage{enumitem} % Includes lists

\frenchspacing % No double spacing between sentences
\linespread{1.2} % Set linespace
\usepackage[a4paper, lmargin=0.1666\paperwidth, rmargin=0.1666\paperwidth, tmargin=0.1111\paperheight, bmargin=0.1111\paperheight]{geometry} %margins
%\usepackage{parskip}

\usepackage[all]{nowidow} % Tries to remove widows
\usepackage[protrusion=true,expansion=true]{microtype} % Improves typography, load after fontpackage is selected

\usepackage{lipsum} % Used for inserting dummy 'Lorem ipsum' text into the template

%--------------------
% Needed by the content below (maths and the results tables)
% -------------------
\usepackage{amsmath,amssymb}
\usepackage{booktabs,array}
\setcounter{secnumdepth}{0} % headings carry their own numbers
\providecommand{\tightlist}{\setlength{\itemsep}{0pt}\setlength{\parskip}{0pt}}

\title{GAP models for the dipole and polarisability of liquid water: what worked and why}
\author{}
\date{}

\begin{document}
\maketitle

\section{1. Scope, metric and protocol}\label{scope-metric-and-protocol}

Every result below is a sparse Gaussian-process (GAP) fit in the gap\_fit estimator: for design
matrix \(\Phi\), targets \(y\), per-frame observation error \(\sigma_\mathrm{obs} = \sigma\sqrt{N}\)
and sparse-sparse prior \(K_{MM}\),

\[
\hat\alpha = \arg\min_\alpha \; \sum_f \frac{\lVert \Phi_f\alpha - y_f\rVert^2}{\sigma_\mathrm{obs}^2} + \alpha^\top K_{MM}\,\alpha ,
\]

solved in whitened coordinates \(\beta = (L^\top\!\otimes I)\alpha\) with \(K_{MM}=LL^\top\), in the
dual (\([DF\times DF]\) Gram, \(D=3\) or \(6\) rows per frame). Fits take one to two minutes, which
is what made the whole study possible: every readout question was settled on the
\emph{representability floor} of the basis before any SNES/Adam training was considered.

Accuracy is quoted as the relative RMSE of the TNEP paper (Xu et al., JCTC 2024),

\[
\mathrm{RRMSE} = \frac{\mathrm{RMSE}}{\mathrm{std}(\text{reference})},\qquad R^2 = 1-\mathrm{RRMSE}^2 ,
\]

computed over components (the three Cartesian dipole components; the three diagonal and the three
off-diagonal Voigt components of \(\alpha\) separately). Note that the code's own \texttt{rrmse}
history divides by the raw target magnitude and is dominated by the 44 \AA{}\(^3\) isotropic trace
of \(\alpha\); it is not comparable to published values.

Protocol: hyperparameters (\(\sigma\), kernel-group amplitudes) were selected on a 400/100 split of
the training set and the winner refitted on all 500 frames before scoring the test set. Where a
table says ``fit 400'', the model was fitted on 400 frames and scored on test directly. Published
reference values: dipole 0.85 \% (same label spread as ours, \(0.20\ e\cdot a_0\) per molecule, so
the comparison is fair); polarisability 16.3 \% (diagonal) / 20.4 \% (off-diagonal), on the
susceptibility \(\chi/\rho\) rather than the cell polarisability, so only the relative numbers are
comparable.

\section{2. The readout family}\label{the-readout-family}

Both properties are read out from per-atom \emph{site functions} multiplying fixed \emph{geometric
kernels} built from the bond vectors \(\mathbf d_{ij} = \mathbf r_i-\mathbf r_j\) of all image pairs
within \(r_c\):

\[
\mu = \sum_i \sum_k s^{(k)}_i\, \mathbf v^{(k)}_i, \qquad
\alpha = \sum_i \sum_k s^{(k)}_i\, \mathbf W^{(k)}_i ,
\]

where \(k\) runs over radial channels \(n\), neighbour species \(t'\) and kernel groups. In the GAP
the site functions are kernel expansions on the SOAP descriptor \(D_i\) of the environment,

\[
s^{(k)}_i = \sum_m \alpha^{(k)}_m\,\big(\hat D_i\cdot\hat D_m\big)^{\xi}, \qquad \xi = 2,
\]

The radial channels are the warped Chebyshev family of the production readout,
\(R_n(r) = T_n\!\big(2(r/r_c-1)^2-1\big)\,f_c(r)\) with a cosine or squared-cosine cutoff,
\(C = 10\) channels, with independent coefficients per (channel, neighbour species) pair.

\section{3. Dipole models}\label{dipole-models}

Three models are documented: the original neighbour-basis GAP (the baseline every change is measured
against), the best \emph{local} model, which keeps the 3 \AA{} readout and adds nothing long-ranged,
and the best model overall, which adds the Ewald field of fixed reference charges as one extra
readout column. All three share the same readout, so switching between them is a change of
configuration, not of code.

\subsection{3.1 The original model (baseline)}\label{the-original-model-baseline}

\[
\mu=\sum_i\sum_{n,t'} s^{(n,t')}_i\sum_{j\in t'}R_n(r_{ij})\,\mathbf d_{ij},
\qquad
s^{(n,t')}_i=\sum_m\alpha^{(n,t')}_m\big(\hat D_i\cdot\hat D_m\big)^2 .
\]

A bond basis of \(C=10\) warped Chebyshev channels per neighbour species on \(r_c=3.0\) \AA{} with a
squared-cosine cutoff, each channel with its own site function: a sparse kernel on the compressed
\(l_{\max}=\alpha_{\max}=4\) SOAP (hard/soft cutoff 3.5/3.0 \AA{}), \(M=250\) farthest-point
environments per species, \(\xi=2\), \(\sigma=10^{-3}\). 10 000 columns. \textbf{Test RRMSE 1.76 \%}
on 500 training frames (1.86 \% on 400); 0.28 ms/atom.

\subsection{3.2 Best local model (no field)}\label{best-local-model-no-field}

Two changes to §3.1, both inside the site function; the readout is untouched.

\begin{enumerate}
\def\labelenumi{\arabic{enumi}.}
\tightlist
\item
  \textbf{A richer descriptor.} \(l_{\max}=\alpha_{\max}=7\) with no compression (the full
  840-dimensional power spectrum instead of the 4-4 compressed one). 1.76 \(\rightarrow\)
  \textbf{1.56 \%} (fit 500; 1.67 on 400). The selected \(\sigma\) stays pinned at the smallest
  value scanned, the signature of a capacity-limited site function, which is why more descriptor
  helped and why \(M=500\) did not (saturated at 250).
\item
  \textbf{A summed two-body site model} (§3.7). The site function becomes
  \(s_i=\sum_m\alpha_m(\hat D_i\!\cdot\!\hat D_m)^2+\delta_{2b}\sum_{m'}\beta_{m'}(\hat{\mathbf q}^{(2)}_i\!\cdot\!\hat{\mathbf q}^{(2)}_{m'})^2\)
  with \(q^{(2)}_{i,nt'}=\sum_{j\in t'}R_n(r_{ij})\) the 20-number radial distribution of the
  environment on the same 3 \AA{}, its own FPS set (\(M=250\)) and its own prior,
  \(\delta_{2b}=0.1\). 1.56 \(\rightarrow\) \textbf{1.49 \%} (fit 500; 1.56 on 400) at
  \(\sigma=3\times10^{-3}\). Columns double to 20 000; the forward pass is unchanged within noise.
\end{enumerate}

\subsection{3.4 Results (test set, 500 training frames unless
stated)}\label{results-test-set-500-training-frames-unless-stated}

RMSE is RRMSE \(\times\) the test-set spread: 6.24 \(e\cdot a_0\) per cell for a dipole component
(0.195 per molecule, 32 molecules).

\begin{center}\footnotesize\setlength{\tabcolsep}{4pt}
\begin{tabular}{@{}>{\raggedright\arraybackslash}p{4.4cm}cccccc@{}}
\toprule
Model & \multicolumn{2}{c}{RRMSE} & \multicolumn{2}{c}{RMSE (fit 500)} & columns & ms/atom \\
\cmidrule(lr){2-3}\cmidrule(lr){4-5}
 & fit 400 & fit 500 & $e\cdot a_0$/cell & $10^{-3}\,e\cdot a_0$/mol & & \\
\midrule
§3.1 original: $l=\alpha=4$ compressed SOAP, bond basis & 1.86 \% & 1.76 \% & 0.110 & 3.4 & 10\,000 & 0.28 \\
$l=\alpha=7$ uncompressed SOAP & 1.67 \% & 1.56 \% & 0.097 & 3.0 & 10\,000 & 0.43 \\
\textbf{§3.2 best local: + 2b site model, $\delta_{2b}=0.1$} & 1.56 \% & \textbf{1.49 \%} & \textbf{0.093} & 2.9 & 20\,000 & 0.45 \\
$l=\alpha=7$ + Ewald field column & 1.17 \% & 1.08 \% & 0.067 & 2.1 & 11\,000 & 1.45 \\
\textbf{§3.3 best: + field + 2b site model, $\delta_{2b}=0.05$--$0.1$} & 1.05 \% & \textbf{0.96--0.97 \%} & \textbf{0.060} & 1.9 & 22\,000 & 1.47 \\
Published TNEP (NEP descriptors, 6 \AA{} radial / 4 \AA{} angular, 200k SNES generations) & & 0.85 \% & 0.053 & 1.7 & & \\
\bottomrule
\end{tabular}
\end{center}

\(R^2=0.9999\) for the best model. The TNEP row's RMSE is the paper's own \(1.7\times10^{-3}\,e\)
per molecule; the label spread is the same as ours, so the two RMSE columns are directly comparable.

\subsection{3.7 The summed 2b site model and its
probes}\label{the-summed-2b-site-model-and-its-probes}

The site function of §3.1 is a single sparse kernel on the SOAP vector. The NEP decomposition of the
site energy by body order suggests the sum

\[
s_i = \sum_m \alpha_m\,(\hat{\mathbf q}_i\cdot\hat{\mathbf q}_m)^{2}
\;+\;\delta_{2b}\sum_{m'} \beta_{m'}\,(\hat{\mathbf q}^{(2)}_i\cdot\hat{\mathbf q}^{(2)}_{m'})^{2},
\qquad
q^{(2)}_{i,\,n t'} = \sum_{j\in t'} R_n(r_{ij}),
\]

a SOAP kernel plus a kernel on the radial distribution of the environment (\(C=10\) warped Chebyshev
channels per neighbour species, 20 numbers), each with its own FPS sparse set (\(M=250\) per
species) and its own \(K_{MM}\) prior, never mixed; \(\delta_{2b}\) is the relative prior amplitude.
Both read out through the same bond basis and field column. A 3b model (angular power spectrum
\(q_{nn'l}\), \(l\le2\)) reads out through an angular kernel
\(\sum_{j,k\ne j}R_aR_b\,P_l(\cos\theta_{jik})\,\mathbf d_{ij}\) instead, so that the bond vector
stays uncontracted.

Val-selected on 400/100 and refit on 500 frames, all with the Ewald field column (test RRMSE): SOAP 1.08 \%; \textbf{SOAP + 2b,
\(\delta_{2b}=0.3\), \(\sigma=3\times10^{-3}\): 1.00 \%}; SOAP + 3b (\(\delta_{3b}=0.3\)) 1.02 \%;
SOAP + 2b + 3b (\(\delta=0.3\)) 1.05 \%. Alone, the 2b model reaches 1.76 \% and the 3b model 2.65
\%, i.e.~neither replaces SOAP, and at \(\delta=1\) every sum is worse than SOAP alone (1.19--1.37
\% on 400 frames against 1.17). The gain is modest (0.08 points, \(\approx\) 3\(\times\) the
\(\sigma\)-to-\(\sigma\) noise) and comes with a shift of the selected \(\sigma\) from \(10^{-3}\)
to \(3\times10^{-3}\): a small-amplitude, smooth two-body component lets the composite tolerate
stronger regularisation, which is where the improvement appears. It costs 2\(\times\) the columns
(22 000) and nothing in fit time. This is the mechanism the TurboGAP 2b \emph{column block} of §7
could not deliver: there the distance features were linear columns inside the one SOAP site
function; here the 2b term is a separately regularised site function of its own.

\textbf{Probe of the 2b knobs and of multiscale} (\texttt{gap\_2b\_probe.py}, same protocol,
forward-pass cost from \texttt{predict\_timed}). The amplitude keeps improving down to
\(\delta_{2b}=0.05\)--\(0.1\) with \(\sigma=1.5\)--\(2\times10^{-3}\) (test 0.96--0.97 \%, val 1.07
against 1.08--1.13 for \(\delta_{2b}=0.1\)--\(0.5\)), then plateaus. Everything that makes the 2b
kernel richer makes it worse: evaluating \(q^{(2)}\) on a longer cutoff than the readout
(\(r_c^{(2)}=4,5,6\) \AA{}: val 1.23, 1.38, 1.49 \% against 1.10 at 3 \AA{}, monotone), 20 channels
(1.29) or 6 (1.21), \(\xi\in\{2,4\}\) (1.18), \(M=500\) (1.10, no change), and the raw polynomial
kernel on the un-normalised distribution, which keeps the neighbour count (val 1.17, test
1.00--1.01). A second SOAP scale (\texttt{gap\_extra\_rcut}) does not help the dipole: alone, 4.5
\AA{} gives 1.18--1.24 on val but 1.21--1.33 on test (against 1.17) and 6.0 \AA{} is clearly worse
(1.43--1.55); stacked on SOAP + 2b it lowers val to 1.04 but the test stays at 1.00 \%. So the 2b
term is what this section said it is --- a weak, smooth, short-range prior --- and the long range of
the dipole is already carried by the field column.

\emph{Cost.} The forward pass (neighbour lists, SOAP, 2b descriptor, geometry and kernel columns;
two processes sharing 16 cores) is 1.45 ms/atom for SOAP alone and 1.47--1.54 ms/atom with the 2b
model (+2--7 \%); a 6 \AA{} 2b cutoff costs 1.6\(\times\), a second SOAP scale at 4.5 \AA{}
1.5\(\times\) and at 6 \AA{} 1.85\(\times\). The Ewald field column is the largest single item: 1.47
ms/atom with it against 0.45 ms/atom without, in a pure-Python Ewald sum.

\textbf{A long readout basis does not transfer.} The published TNEP reads the dipole out through
\(\partial U_i/\partial\mathbf r_{ij}\) of a learned 6 \AA{} two-body basis, i.e.~a bond basis with
learned radial functions reaching 6 \AA{} and site functions tied to one network. The kernel
analogue --- the same SOAP + 2b site functions reading out through a bond basis with \(r_c=4.5\) or
6 \AA{}, 2b descriptor kept at 3 \AA{} --- is monotonically worse: 0.97 \(\rightarrow\) 1.11
\(\rightarrow\) 1.41 \% at 10 channels, 1.46 \% at 20, while the training RRMSE falls from 0.60 to
0.03 \%. Without the field column the 6 \AA{} readout gives 1.42 \% against 1.49 \% at 3 \AA{}, so
the field's whole gain disappears once the basis can reach 6 \AA{}: the long bond basis fits what
the field explains on the training set and does not transfer it. With one independent site function
per (channel, species) the long basis simply has too much freedom for 500 frames; the TNEP can
afford the range because every channel is the derivative of a single 1097-parameter scalar field. In
a kernel model the long range has to enter as a fixed physics column with one site function, which
is what the field column is.

\section{4. Polarisability models}\label{polarisability-models}

\subsection{4.1 Form}\label{form}

The \textbf{original model} is the first two terms below --- the isotropic and traceless bond
groups, \(C=10\) channels each on \(r_c=4\) \AA{} with a cosine cutoff, site functions on the
\(l=\alpha=7\) compressed SOAP, \(M=250\), \(\sigma=0.1\), 20 000 columns: \textbf{16.6 / 19.9 \%}
diagonal / off-diagonal on 500 training frames (16.9 / 20.2 on 400), 0.63 ms/atom. Every later
number is measured against it. The full model has four even kernel groups, each with its own site
function:

\begin{align*}
\alpha ={}& \sum_i\sum_{n,t'}\Big[
 s^{(n,t')}_i \sum_{j\in t'} R_n(r_{ij})\,|\mathbf d_{ij}|^2\,\mathbf I
\;+\;
 u^{(n,t')}_i \sum_{j\in t'} R_n(r_{ij})\Big(\mathbf d_{ij}\mathbf d_{ij}^\top - \tfrac{|\mathbf d_{ij}|^2}{3}\mathbf I\Big)\Big] \\
&+\;\sum_i\sum_{a\le b} w^{(ab)}_i\,\mathrm{sym}\!\big(\mathbf v^{(a)}_i\otimes\mathbf v^{(b)}_i\big)_{\text{traceless}}
\;+\;
\sum_i z_i\,\Big[\sum_j \big((\mathbf I - \mathbf A\mathbf T)^{-1}\mathbf A\big)_{ij}\Big]_{\text{traceless}} .
\end{align*}

The first two are the isotropic and traceless \emph{bond} groups (\(C=10\) channels each,
\(r_c = 4\) \AA{}, cosine cutoff); the third is the \emph{cross} group built from four radial
channels of the dipole kernel vector \(\mathbf v^{(a)}_i=\sum_j R_a(r_{ij})\mathbf d_{ij}\); the
fourth is the traceless part of the symmetrised Applequist induced-coupling site tensor
\(\mathbf W_i=\tfrac12(\mathbf R_i+\mathbf R_i^\top)\),
\(\mathbf R_i=\sum_j\big[(\mathbf I-\mathbf A\mathbf T)^{-1}\mathbf A\big]_{ij}\), with fixed
\(\alpha_\mathrm O=1.45\) \AA{}\(^3\), \(\alpha_\mathrm H=0\) and the Ewald dipole--dipole tensor
\(\mathbf T\) (intramolecular pairs excluded); it is one extra column per site function, scaled once
to the RMS of the bond columns (production: \texttt{pol\_field\_groups={[}"induced"{]}},
\texttt{pol\_induced\_alpha}, pinned \texttt{pol\_field\_scale}). Site functions on the
\(l=\alpha=7\) compressed SOAP, \(M=250\), \(\sigma=0.1\), 38 500 columns. This is the recommended
model: it adds a single physics column to the base and nothing else, costs no forward time, and
takes 1.9 points off the diagonal and 1.4 off the off-diagonal (§4.2).

\subsection{4.2 Results (test set, diagonal / off-diagonal; all fits on 500 frames unless
stated)}\label{results-test-set-diagonal-off-diagonal-all-fits-on-500-frames-unless-stated}

RMSE is RRMSE \(\times\) the test-set spread of the cell tensor's components: 0.887 \AA{}\(^3\) for
a diagonal component (about the isotropic \(44.4\pm0.7\) \AA{}\(^3\)), 0.468 \AA{}\(^3\) for an
off-diagonal one; per molecule divide by 32.

\begin{center}\footnotesize\setlength{\tabcolsep}{4pt}
\begin{tabular}{@{}>{\raggedright\arraybackslash}p{5.0cm}ccccc@{}}
\toprule
Model & \multicolumn{2}{c}{RRMSE} & \multicolumn{2}{c}{RMSE (\AA{}$^3$/cell)} & ms/atom \\
\cmidrule(lr){2-3}\cmidrule(lr){4-5}
 & diag & off-diag & diag & off-diag & \\
\midrule
Predict the component means & 100 \% & 100 \% & 0.887 & 0.468 & \\
Single $\mathbf d\otimes\mathbf d$ kernel (SNES, 5900 generations) & \multicolumn{2}{c}{mean-predictor floor} & & & \\
Two bond groups, SNES-trained ANN (15 neurons, 25\,k parameters, $\sigma_0=0.02$, 20\,000 generations) & 35.7 \% & 27.5 \% & 0.317 & 0.129 & \\
\textbf{Two bond groups (iso + traceless): the original model, §4.1} & 16.6 \% & 19.9 \% & 0.147 & 0.093 & 0.63 \\
+ cross group, 4 channels (the base) & 14.9 \% & 16.0 \% & 0.132 & 0.075 & 0.86 \\
\textbf{+ traceless induced-coupling group (§4.1, recommended)} & \textbf{14.0 \%} & \textbf{14.6 \%} & \textbf{0.124} & \textbf{0.068} & 0.80 \\
+ induced coupling + field quadrupole (§4.5) & 13.9 \% & 14.5 \% & 0.123 & 0.068 & 0.82 \\
+ all five long-range groups (§4.5) & 13.9 \% & 14.3 \% & 0.123 & 0.067 & 0.84 \\
+ second SOAP scale at 6.5 \AA{}, $\sigma=0.3$ (§4.4) & 13.9 \% & 16.2 \% & 0.123 & 0.076 & 2.5 \\
+ induced coupling + quadrupole + 6.5 \AA{} scale, $\sigma=0.3$ & 13.0 \% & 15.2 \% & 0.115 & 0.071 & 2.6 \\
\quad same, $\sigma=0.1$ & 13.9 \% & 14.6 \% & 0.123 & 0.068 & 2.6 \\
Published TNEP (susceptibility $\chi/\rho$ per molecule; its own RMSE $54\times10^{-3}$ / $37\times10^{-3}$ bohr$^3$) & 16.3 \% & 20.4 \% & & & \\
\bottomrule
\end{tabular}
\end{center}

Forward-pass cost from \texttt{predict\_timed} on the 96-atom test frames, one process on 12
threads; per-frame Python overhead dominates at this size, so only the ratios are meaningful.

The SNES row is the same two-group basis as the original GAP model, trained as a network instead of
solved as a kernel: the final-generation checkpoint
(\path{models/n15_q840_pop100_20261018_231703}) scored on the test set gives a pooled RMSE of
0.2395 \AA{}\(^3\) and per-component \(R^2\) of 0.87--0.88 (diagonal) and 0.91--0.93 (off-diagonal),
i.e.~about twice the diagonal floor and 1.4\(\times\) the off-diagonal floor of its own basis after
20 000 generations. This is the optimiser gap of §8, not a basis limit.

\subsection{4.3 Why it works}\label{why-it-works}

\textbf{Scale separation.} A single \(w\,\mathbf d\otimes\mathbf d\) bond term is purely
longitudinal: its trace \(w|\mathbf d|^2\) and its anisotropy
\(w(\mathbf d\mathbf d^\top-|\mathbf d|^2\mathbf I/3)\) are locked to one coefficient, while a real
bond has independent \(\alpha_\parallel\) and \(\alpha_\perp\) and the isotropic (\(44\pm0.7\)
\AA{}\(^3\)) and anisotropic (\(\pm0.5\) \AA{}\(^3\)) parts of the cell tensor differ by two orders
of magnitude. Splitting the groups is the classical bond-polarisability decomposition and mirrors
GPUMD's scalar-plus-virial pair. In the ANN the split was decisive (mean-predictor floor to 0.25
\AA{}\(^3\)); in the GAP the per-channel prior already absorbs scale, so group amplitudes and the
shear weight moved results by half a point.

\textbf{Two-neighbour terms.} Expanding \(\mathbf v\otimes\mathbf v\) gives
\(\sum_{j,k}R_aR_b\,\mathbf d_{ij}\otimes\mathbf d_{ik}\), the bond-angle anisotropy that GPUMD
obtains from the derivative of its angular descriptors and that bond-axis tensors cannot represent.
It saturated at four radial channels (6, 8 and 10 channels were flat or worse), so the angular
information is low-dimensional.

\textbf{Long-range induced coupling.} The learning curve (24.4 / 20.9 \% at 100 frames, 18.6 / 17.8
at 200, 15.3 / 16.3 at 400, 14.9 / 16.0 at 500; diagonal \(\propto N^{-0.33}\)) showed the
off-diagonal flattening first, which pointed at a form limit for the anisotropy rather than a data
limit. The induced dipole--dipole coupling \(\alpha_i\mathbf T_{ij}\alpha_j\) is traceless, decays
as \(r^{-3}\) and reaches well beyond 4 \AA{}. Solving it as a linear response with fixed
\(\alpha_\mathrm O\) (one \(3N\times3N\) solve per frame) and keeping only the traceless part of the
per-site response improved both the diagonal and the off-diagonal; the full tensor's isotropic part
competed with the isotropic bond group and made the diagonal worse (16.7 \%). The result was
identical for \(\alpha_\mathrm O = 1.0\) and \(1.45\) \AA{}\(^3\), i.e.~the geometry of the coupling
matters and its amplitude is absorbed by the site function.

\subsection{4.4 A second SOAP scale}\label{a-second-soap-scale}

The one addition that moved the polarisability in the 2026-10-20 probes is multiscale: a second SOAP
site function on the same \(l=\alpha=7\) compressed descriptor built with a longer hard cutoff
(\texttt{gap\_extra\_rcut}), its own FPS set and prior, reading out through the same three groups.
Val-selected and refit on 500 frames: 5.5 \AA{} 14.2 / 16.1 \%, \textbf{6.5 \AA{} 13.9 / 16.2 \%},
7.5 \AA{} 14.0 / 16.2 \% (SOAP alone 14.9 / 16.0). The gain is on the diagonal only and saturates at
6.5 \AA{}; the selected \(\sigma\) moves from 0.1 to 0.3; \(\sigma=0.5\) is worse (14.1 / 17.2).
This is consistent with §4.3: the isotropic polarisability of a molecule depends on its wider
environment (the learning curve said the diagonal was data-limited, i.e.~there was information the 4
\AA{} SOAP could not reach), whereas the anisotropy is a form limit that neither more data nor a
longer descriptor touches --- the induced-coupling group of §4.1 remains the only thing that has
moved the off-diagonal. The 2b site model, by contrast, hurts \(\alpha\) at every amplitude and
range tried (\(\delta_{2b}=0.1\): 16.2 / 17.8 val; \(\delta_{2b}=0.3\) on a 6 \AA{} cutoff: 16.0 /
18.0; stacked on the 5.5 \AA{} scale: 16.1 / 17.6).

\end{document}
